%% LyX 2.4.4 created this file.  For more info, see https://www.lyx.org/.
%% Do not edit unless you really know what you are doing.
\documentclass[10pt,english]{article}
\usepackage[T1]{fontenc}
\usepackage[latin9]{inputenc}
\usepackage{amssymb}
\usepackage{geometry}
\geometry{verbose,tmargin=1in,bmargin=1in,lmargin=1in,rmargin=1in}

\makeatletter

%%%%%%%%%%%%%%%%%%%%%%%%%%%%%% LyX specific LaTeX commands.
%% Because html converters don't know tabularnewline
\providecommand{\tabularnewline}{\\}

%%%%%%%%%%%%%%%%%%%%%%%%%%%%%% User specified LaTeX commands.
\usepackage{hyperref}
\usepackage[usenames,dvipsnames]{xcolor} %adds additional color options
\hypersetup{colorlinks=true, urlcolor=Blue}

\usepackage{sectsty} %%one way to change section header font sizes.
\sectionfont{\large}

\usepackage{titlesec}
\titlespacing*{\section}{0pt}{10pt}{0pt}

\usepackage{multicol}
\usepackage{datetime}

\longdate

\def\formatdateny{\csname noyear\languagename\expandafter\endcsname\formatdate}
\def\noyearenglish#1, \the\@year{#1}
\let\noyearamerican\noyearenglish
\let\noyearbritish\noyearenglish
\def\noyearfrench#1\space\number\@year{#1}
\let\noyeargerman\noyearfrench
\let\noyearaustrian\noyeargerman
\let\noyearswedish\noyearfrench
\let\noyearbreton\noyearfrench
\def\noyearrussian#1\ \number\@year~\cyrg.{#1}
\def\noyearspanish#1\ de~\number\@year{#1}
\let\noyearcatalan\noyearspanish
\def\noyearbasque#1\number\@year.eko\space{#1}
% Added by lyx2lyx
\setlength{\parskip}{\medskipamount}
\setlength{\parindent}{0pt}

\makeatother

\usepackage{babel}
\begin{document}
\begin{center}
{\Large\textbf{Johnson C. Smith University}}\textbf{}\\
\textbf{Introductory Linear Algebra}\\
\textbf{MTH 336(A)~~|~~9AM-9:50AM SHB 303~~|~~Fall 2013}\\
\textbf{\rule{1\columnwidth}{1pt}}
\par\end{center}

\vspace{-.3cm}\begin{multicols}{2}

\section*{Instructor Information}

\textbf{Name: }Dr.\,James Ashe\\
\textbf{Office location: }SHA 307 Davis Hall (back)\\
\textbf{Email: }\href{mailto:jashe@jcsu.edu}{jashe@jcsu.edu}\\
\textbf{Office phone: }(704)\,378-1037\\
\textbf{Office hours:}

\scriptsize%
\begin{tabular}{lll}
10:00 AM--11:00 AM & MWF & \tabularnewline
12:00 PM--1:00 PM & MW & \tabularnewline
12:00 PM--1:00 PM & F & (by appointment only)\tabularnewline
1:00 PM--2:00 PM & MW & (by appointment only)\tabularnewline
11:00 AM--12:00 PM & TR  & (by appointment only)\tabularnewline
\end{tabular}\small

\section*{Course Description}

This is a \textbf{3 credit hour}\textbf{\emph{ }}course; Systems of
linear equations, vector spaces, and linear transformations, plus
theory and applications of matrices and determinants. Prerequisite:
MTH 231 or consent of Department. Prerequisite: MTH 231 (Calculus
I). 

Instructional methods will primarily consist of lectures and in-class
discussions, however other methods such collaborative projects and
student teaching may be adopted according to the particular needs
of the class.\end{multicols}

\section*{Course Objectives}

Students should be able to demonstrate competency in the techniques
and topics outlined in the section above. Additionally students should
develop general problem solving skills and apply tools acquired in
the course to interpret and solve problems.

\section*{Course Materials }

\textbf{Textbook:}\emph{ \href{http://www.amazon.com/Linear-Algebra-A-Geometric-Approach/dp/1429215216/ref=pd_sim_sbs_b_1}{Linear Algebra: A Geometric Approach, Second Edition}}

\textbf{Calculator (recommended): }TI-83, TI-83+, or TI-84+. Any other
calculator model must be approved by the instructor.\\
Students should bring their calculator and textbook to class.

\section*{Grades}

Grades will be taken from homework, quizzes, and Exams as follows:\textbf{}\\
\textbf{Grading scale: $\mbox{A}\geq90>\mbox{B}\geq80>\mbox{C}\geq70>\mbox{D}\geq60>\mbox{F}$}\\
\textbf{Evaluation: }Quizzes$=20\%$; Exams 1, 2, and 3 $=20\%$ each;
Final Exam $=20\%$.

\textbf{Quizzes: }Quizzes will consist of any graded assignment which
is not an exam. This will include traditional quizzes, take home assignments,
and in-class projects. The lowest quiz grade will be dropped. \textbf{}\\
\textbf{Final Exam: }The final Exam will be a comprehensive in-class
exam; see the exam schedule \\
\href{http://www.jcsu.edu/happenings/exam-schedule}{jcsu.edu/happ enings/exam-schedule}.
If beneficial, the final exam score will replace the lowest exam score.

\section*{Class Policies}

JCSU policies regarding attendance, academic integrity, grades, and
change of enrollment will be enforced. It is the student\textquoteright s
responsibility to be aware of these policies and of related deadlines
as well as any announcements or syllabus adjustments (written or oral)
made in class. Please see the student handbook for details, \href{http://www.jcsu.edu/docs/handbook.pdf}{jcsu.edu/docs/handbook.pdf}.

\textbf{Attendance: }Class attendance is required for all JCSU students.
Each student is allowed as many hours of absence per term as credit
hours(s) received (not to exceed 4) for the class. Students who exceed
the maximum number of absences may receive a failing grade for the
course. 

Students who may miss classes while representing the University in
an official capacity are exempt from regulations governing absences.
However, absence from class for official University business does
not relieve the student from responsibility for any class assignments
that may be missed during the period of absence.

\textbf{Missed quizzes: }There will be no make-ups for in-class quizzes
unless the class was missed for a University sponsored event, and
advance notice of the absence was given by the student. In this case
the assignment will be excused or a make-up will be given depending
on circumstances. Late take-home quizzes will have 25\% deducted for
each day it is late. \\
\textbf{Missed exams: }If a student misses an exam for a verifiable
excuse, they should send me an e-mail \emph{as soon as possible }so
that a make-up can be arranged. Make-up exams may be significantly
more difficult. If you will be gone for any reason, it will be \emph{your
}responsibility to inform me at least a week in advance to schedule
to take the exam early. 

\textbf{Student support services: }If you need course adaptations
or accommodations for a documented disability please contact the Office
of Disability Services. \\
\href{http://www.jcsu.edu/directory/federal-requirements/general-institutional-information/jcsu-disability-services}{jcsu.edu/directory/federal-requirements/general-institutional-information/jcsu-disability-services}
,\\
phone: (704)\,398-1116.

The Office of Counseling Services provides confidential intakes and
assessments, personal counseling and referrals to appropriate resources
to all enrolled students on a walk-in or appointment basis. \\
\href{http://www.jcsu.edu/student/office-of-counseling-services}{www.jcsu.edu/student/office-of-counseling-services},
\\
phone (704)\,378-1044.

\section*{Course Content}

This is an outline of topics that will be covered in this course.
Additional sections may be covered as time allows

\begin{multicols}{2}

\textbf{Chapter 1: Vectors and Matrices}\\
1. Vectors\\
2. Dot Product\\
3. Hyperplanes in $\mathbb{R}^{n}$\\
4. Systems of Linear Equations\\
5. Theory of Linear Systems\\
6. Applications (if time allows)

\textbf{Chapter 2: Matrix Algebra}\\
1. Matrix Operations\\
2. Linear Transformations\\
3. Inverse Matrices\\
4. Elementary Matrices\\
5. The Transpose

\textbf{Chapter 3: Vector Spaces}\\
1. Subspaces of $\mathbb{R}^{n}$\\
2. The Four Fundamental Subspaces\\
3. Linear Independence and Basis\\
4. Dimension and Its Consequences

\textbf{Chapter 4: Projections and Linear Transformations}\\
1. Inconsistent Systems and Projections\\
2. Orthogonal Bases\\
3. The Matrix of a Linear Transformation

\textbf{Chapter 5: Determinants}\\
1. Properties if Determinants\\
2. Orthogonal Bases

\textbf{Chapter 6: Eigenvalues and Eigenvectors}\\
1. The Characteristic Polynomial\\
2. Diagonalizability\\
4. The Spectral Theorem 

\end{multicols}

\section*{Important Dates}

\noindent\begin{minipage}[t]{1\columnwidth}%
\begin{tabular}{rcl}
\formatdateny{11}{9}{2013} &  & Last day to Add/Drop \tabularnewline
\formatdateny{18}{9}{2013} &  & Exam 1\tabularnewline
\formatdateny{23}{10}{2013} &  & Exam 2\tabularnewline
\formatdateny{27}{11}{2013} &  & Exam 3\tabularnewline
TBA &  & Final Exam\tabularnewline
 &  & \tabularnewline
\end{tabular}%
\end{minipage}

With exception of the final, the above exam dates are subject to change.
JCSU's Fall 2012 academic calendar can be found here: \href{http://www.jcsu.edu/happenings/academic-calendar}{jcsu.edu/happenings/academic-calendar},
and the exam schedule can be found here: \href{http://www.jcsu.edu/happenings/exam-schedule}{jcsu.edu/happenings/exam-schedule}.

While I do not foresee any significant changes to the items in this
syllabus during the semester, I reserve the right to make appropriate
changes. Students will be notified of any such changes in writing.
\end{document}
